\begin{statement}{C.2.CR.6}{Proved}{The elliptic contraction has a planar conductor and only double transverse branches at its image}
In the remaining case $g=1$ of C.2.CR.1, let
$p=\pi(C)$ and $q=\eta(p)$. The conductor-image
scheme is a local complete intersection at $q$.
More precisely, ambient coordinates can be chosen
so that its ideal is
\[
                         (x,g(y,z)).               \tag{C.CR.21}
\]
Every global conductor component passing through
$q$ has generic transverse type $A_k$ for some
$k\ge1$. For type $A_{2s-1}$, its two transverse
normalization branches have conductor coefficient
$s$ each; their global covering degrees sum to two.
For type $A_{2s}$ there is one coefficient-$2s$
normalization branch, of covering degree one.
\end{statement}
\begin{proof}
The residue comparison (C.CR.12) is valid also
when $g=1$. By (C.CR.13) its nonzero one-dimensional
quotient requires an element of the conductor
summand $I_p$ with nonzero ambient linear part.
Call this element $x$ and extend it to coordinates
$(x,y,z)$ on the smooth threefold.

The full conductor ideal contains $x$ and has a
pure one-dimensional Cohen--Macaulay quotient.
After quotienting the ambient local ring by $x$,
its remaining ideal is an unmixed height-one
ideal in the regular two-dimensional ring
$\mathbb C\{y,z\}$. Such an ideal is principal:
it is the intersection of its height-one primary
localizations, and each is a power of a principal
prime in this factorial ring. Equivalently its
quotient has depth one, so the ideal is a rank-one
reflexive module on this smooth surface and
is invertible. This proves (C.CR.21), retaining
all multiplicities in $g$.

We give the precise adjunction comparison behind
the transverse multiplicity assertion. For a
reduced plane curve germ $D\subset W$, let
$\nu:\widetilde D\to D$ be its normalization,
and let $\rho:\widetilde W\to W$ be an embedded
resolution with smooth disjoint strict transform
$\widehat D$. Finite curve duality identifies
\[
 \mathfrak c_D\,\omega_D=\nu_*\omega_{\widetilde D}.
\]
The residue sequence on $\widetilde W$ and
Grauert--Riemenschneider vanishing identify the
right side with the residue of
$\rho_*\omega_{\widetilde W}(\widehat D)$.
In a local generator $du\wedge dv/f_D$ of the
ambient meromorphic residue description, a
conductor element $a$ therefore lifts to a form
\[
                     a\,du\wedge dv/f_D
\]
whose pullback has no exceptional pole. On the
first plane blowup, with exceptional curve $E$,
the canonical divisor has coefficient one and
the total transform of $D$ has coefficient
$m=\operatorname{mult}_0 D$. The coefficient
of this pulled-back form along $E$ is
$\operatorname{ord}_0(a)+1-m$. Hence
\[
 \mathfrak c_D\subset\mathfrak m_D^{\,m-1}
 \quad\hbox{in the sense that every ambient
 representative has order at least }m-1.
 \tag{C.CR.22}
\]
Adding a multiple of $f_D$ does not change this
order bound, since $f_D$ has order $m$.

At a general smooth point of any conductor
component through $q$, the ambient coordinate
$x$ remains in its conductor ideal and has
a nonzero transverse linear part: its zero
set is a smooth ambient plane containing the
conductor curve. Applying (C.CR.22) to a
transverse plane slice gives $m\le2$.
The slice is singular, so $m=2$.
Weierstrass preparation, completion of the
square, and absorption of a unit then put
a reduced multiplicity-two plane curve in
the form $u^2-v^{k+1}=0$.

For $k=2s-1$, the two branches are
$u=\pm v^s$; their intersection multiplicity
is $s$, and the conductor on each normalized
branch is $(v^s)$. For $k=2s$, the normalization
is $v=t^2$, $u=t^{2s+1}$, whose semigroup
has conductor $2s$, giving $(t^{2s})$.
These calculations specify both the coefficients
and the numbers of transverse branches; finite
normalization over the global image component
expresses those numbers as the asserted sums
of covering degrees.
\end{proof}


% C.2.UN.36 from part_C2_core_lattices.tex

